\documentclass[lang=cn,12pt,scheme=chinese,blue,oneside]{elegantbook}

\usepackage{xeCJK}
\setCJKmainfont{SimSun}
\setCJKsansfont{SimSun}
\setCJKmonofont{SimSun}
\setmainfont{Times New Roman}
\setsansfont{Arial}

\renewcommand{\baselinestretch}{1.35}

\usepackage{geometry}
\geometry{
  a4paper,
  left=24mm,
  right=24mm,
  top=24mm,
  bottom=24mm,
  headheight=2.1cm,
  headsep=4mm,
  footskip=10mm
}

\usepackage{graphicx}
\usepackage{booktabs}
\usepackage{tabularx}
\usepackage{enumitem}

\setlist[itemize]{itemsep=0.4em, topsep=0.6em}
\newcommand{\ProjectCoverImage}{assets/cover.png}
\newcommand{\ShowcaseCover}{%
  \clearpage
  \newgeometry{margin=0pt}%
  \begin{titlepage}
    \thispagestyle{empty}%
    \noindent\includegraphics[width=\paperwidth,height=\paperheight]{\ProjectCoverImage}%
  \end{titlepage}
  \restoregeometry
  \clearpage
}


\title{星语灵境}
\subtitle{面向 ASD 儿童的情绪训练与社交干预系统商业计划书}
\author{星语灵境项目组}
\institute{挑战杯参赛方案}
\date{2026年5月}
\version{feature/profile-edit 分支修订版}
\bioinfo{项目类型}{科技向善·特殊儿童教育支持}

\definecolor{structurecolor}{RGB}{212,154,58}
\definecolor{main}{RGB}{212,154,58}
\definecolor{second}{RGB}{44,62,80}
\definecolor{third}{RGB}{244,238,221}
\colorlet{coverlinecolor}{second}

\setcounter{tocdepth}{2}

\begin{document}
\frontmatter
\maketitle

\clearpage
\thispagestyle{empty}
\begin{center}
\vspace*{2.3cm}
{\Huge\color{structurecolor}\textbf{读懂沉默，回应星光}}

\vspace{1.1cm}
\begin{minipage}{0.84\textwidth}
\begin{center}
\large
\centering
{\kaishu 我们不把它定义为一个冷冰冰的训练工具，}\\[0.35em]
{\kaishu 而是定义为一个能陪孩子慢慢练习社交表达的数字伙伴。}

\vspace{0.8cm}
面向自闭症谱系障碍（ASD）儿童，\\
围绕情绪识别、互动训练、过程记录与反馈评估，\\
构建一套\textbf{可演示、可落地、可扩展}的训练系统原型。
\end{center}
\end{minipage}

\vfill
\PlaceholderFigure
  {封面插图占位}
  {建议后续替换为封面主视觉插图。}
\vfill
\end{center}
\clearpage

\tableofcontents
\mainmatter

\chapter{执行摘要}

\section{项目定位}
“星语灵境”是一套面向自闭症谱系障碍（ASD）儿童的情绪训练与社交支持系统。它并不试图用技术替代家庭、治疗师或教师的工作，而是希望在孩子、家庭与机构之间，补上一层更稳定、更温和、也更容易持续使用的数字化支持。围绕情绪识别、情境理解、自我表达与训练记录，系统把一次次零散练习组织成一条可回看、可沟通、可逐步优化的成长路径。

从商业角度看，本项目并不是面向泛儿童市场的通用应用，而是聚焦于 ASD 儿童训练场景中的真实需求。它既服务于儿童训练本身，也服务于家长、治疗师和康复机构对过程记录、阶段反馈和训练协同的需求。因此，星语灵境的价值不只在“让孩子完成一题”，更在“让支持者看见一次练习究竟发生了什么”。

\section{当前基础}
当前分支已经形成可运行的系统原型。系统支持儿童档案创建与管理，支持五类训练活动的进入与作答，支持基于题库的发题机制、答题记录、训练总结与治疗师看板，也已经打通了基于摄像头的实时面部情绪识别链路。围绕训练过程，系统能够根据儿童作答表现和情绪状态，给出提示、调节难度，并在训练结束后形成汇总反馈。

这意味着本项目已经不是停留在概念描述上的想法，而是具备了演示、试点讨论和继续迭代的原型基础。当前阶段已经打通的能力，应当作为计划书中的事实基础；尚处于规划阶段的能力，则应在后文中作为后续建设方向审慎呈现。

\section{核心判断}
围绕 ASD 儿童情绪训练与社交支持，最突出的矛盾并不是“内容不够丰富”，而是优质训练资源供给不足、训练结果难以迁移到真实生活，以及家庭与机构之间缺少连续、清晰、可信的过程记录。星语灵境的核心判断是：真正有价值的产品，不是简单增加刺激和功能，而是降低训练启动成本，改善训练连续性，并帮助不同支持者形成更一致的观察与判断。

基于这一判断，本项目的落地路径将优先从机构和专业场景切入，通过试点合作验证产品可用性与适配度，再逐步完善产品标准化能力，延展到家庭陪伴和更广泛的教育支持场景。技术在这里承担的是支撑角色，而不是替代角色。

\PlaceholderFigure
  {执行摘要主视觉占位}
  {建议替换为“儿童训练 + 家庭陪伴 + 数据反馈”三部分统一主视觉。}

\chapter{项目缘起与真实需求}

\section{孩子与支持者的共同需要}
对许多 ASD 儿童而言，理解情绪线索、进入社交情境和表达自身感受，往往需要更长时间、更稳定节奏和更细致支持。与此同时，家长、治疗师、特教老师和康复机构也在长期承担陪伴与干预任务。他们真正面对的，不只是“孩子需要练习”，更是“如何让练习能够持续、可理解、可协同”的问题。

这使得 ASD 儿童支持场景天然具有三个特点：训练周期长，单次进步细小；对陪伴方式和刺激强度要求更高；家庭与机构之间需要频繁沟通、共同判断。也正因为如此，这一领域对数字化工具的期待，从来不是一款普通儿童应用，而是一种更贴近训练节奏、更尊重个体差异的支持工具。

\section{训练现场的真实难题}
当前训练场景中的困难，主要集中在三个方面。第一，优质干预资源高度依赖人工投入，机构师资紧张、家庭陪伴成本高，很多训练很难长期保持稳定节奏。第二，传统练习材料多以卡片、图片和标准化课堂为主，孩子即使在训练材料中学会识别，也未必能够顺利迁移到真实生活情境。第三，训练过程往往缺少连续记录，家庭看到的是日常状态，机构看到的是训练表现，但这些观察很难被组织成可共同阅读的过程信息。

对家长来说，最常见的困惑是“孩子今天到底在哪些地方有变化”；对治疗师来说，最常见的压力是“如何在有限时间里兼顾训练、观察和记录”；对机构来说，最难解决的是“怎样把经验沉淀成可复用、可扩展的服务能力”。这些问题共同指向了一个需求，即训练工具不仅要能用，还要能帮助不同支持者形成更清晰的共识。

\section{我们想补上的那一层支持}
星语灵境的切入点，不是把训练包装成更花哨的娱乐内容，而是围绕 ASD 儿童的训练节奏，设计一套更容易进入、也更容易回看的数字化练习流程。儿童端通过低刺激、陪伴式的互动活动降低进入门槛；支持者端通过训练总结和数据看板提高沟通效率；而系统中的状态感知与难度调节机制，则用于帮助训练更贴近儿童当下状态。

这一本质上是一种“训练支持系统”而不是“单次闯关产品”。它服务的不是一次表现，而是一段持续练习；它关心的不是单个对错，而是过程中的理解、稳定性和变化趋势。对于 ASD 场景而言，这种定位比单纯追求功能数量更重要。

\PlaceholderFigure
  {场景痛点图占位}
  {建议替换为“家庭陪伴、机构训练、过程记录”三类痛点并列示意图。}

\chapter{市场机会与竞争判断}

\section{谁最需要这样的产品}
星语灵境面对的是一个多角色共同参与的使用场景。儿童是直接使用者，他们进入训练活动，与机器人伙伴、题目内容和反馈机制互动；家长和康复机构是核心决策者，他们决定是否采用、是否持续使用以及是否愿意投入时间和成本；治疗师、特教老师和学校则是重要协同方，他们决定产品能否真正嵌入训练流程，形成稳定使用习惯。

因此，本项目的市场切入不应从大众儿童应用逻辑出发，而应从专业训练支持场景出发。对首批用户而言，产品是否尊重孩子节奏、是否减轻记录负担、是否能支撑家庭与机构沟通，比功能是否“热闹”更重要。

\section{这件事为什么值得做}
当前市场上并不缺少儿童应用，也不缺少情绪识别、互动游戏或陪伴式内容产品，但真正贴近 ASD 儿童情绪训练场景、同时兼顾陪伴感、训练逻辑和过程记录能力的产品仍然相对稀缺。大量大众产品更擅长提高易用性，却难以回应康复与特教场景中的细致需求；不少研究原型在方法上有启发，但距离稳定落地和持续使用仍有距离。

这意味着市场机会不在于与通用教育产品拼规模，而在于为一个长期存在、却长期缺乏合适工具的专业场景，提供更准确的解决方案。对于康复机构、公益项目、特教支持场景和部分家庭陪伴场景而言，只要产品能够真正降低训练门槛、提升反馈质量，就具备被持续采用的基础。

\section{现有方案与我们的不同}
如果把现有相关方案进行归类，大体可以看到三条路径。第一类是研究型原型，往往有较强的方法论和实验背景，但在产品稳定性、内容组织和持续运营上仍需完善。第二类是面向大众儿童的情绪或教育应用，这类产品更容易上手，但对 ASD 场景的适配普遍不够深入。第三类是机构辅助工具，它们更接近专业使用需求，但未必在儿童交互体验、陪伴表达和训练连续性上形成完整闭环。

星语灵境当前的竞争重点，不是做“最多功能”的产品，而是做“更合适进入专业训练场景”的产品。它将儿童端活动、过程记录、总结反馈和支持者看板串联为一个闭环，并以可解释的掌握度追踪和规则自适应机制支撑训练节奏，这构成了与通用产品之间最核心的差异。

\section{星语灵境的定位}
本项目的差异化定位可以概括为三点。第一，强调陪伴式训练，而不是刺激堆叠式训练；第二，强调过程可读，而不是只给出单次结果；第三，强调逐步落地，而不是超前承诺。对于 ASD 儿童支持场景而言，这样的定位虽然不夸张，却更容易建立长期信任。

\PlaceholderFigure
  {竞争定位图占位}
  {建议替换为“专业度、陪伴感、过程记录能力”三维定位图。}

\chapter{产品方案与现阶段基础}

\section{产品想提供怎样的陪伴}
星语灵境围绕儿童训练与支持者协同，形成了一个由儿童端、记录反馈端和状态感知能力共同组成的产品方案。儿童端承担练习发生的主要空间，孩子在其中完成情绪识别、情境理解和表达练习；记录反馈端承担训练沉淀与沟通功能，帮助家长和治疗师回看训练结果；状态感知能力则用于在训练过程中识别儿童当下表现，并适度触发提示或节奏调整。

在产品气质上，系统强调低刺激、陪伴式和渐进式。机器人伙伴、视觉表达和训练反馈都尽量避免过强压迫感，而是让训练更接近被耐心接住的过程。对于 ASD 儿童而言，能否愿意进入、能否在较低负担下保持参与，往往比一次训练中增加多少“高级功能”更关键。

\section{五类训练活动如何展开}
当前原型中的五类训练活动并不是彼此割裂的模块，而是围绕“识别情绪、理解情境、表达自己”逐步展开。情绪大侦探和表情连连看承担基础识别功能，帮助孩子在相对明确的线索中建立情绪概念；拼脸大师进一步引导孩子理解表情构成，帮助他们从记忆答案走向观察特征；社交小剧场把情绪放入具体情境之中，让孩子开始理解人物关系、事件变化与情绪反应之间的联系；心情日记则把训练从“理解别人”延伸到“表达自己”，让孩子的感受被看见、被记录、也被尊重。

这样的活动安排，一方面照顾了训练难度的递进，另一方面也照顾了 ASD 儿童在进入真实社交场景之前，往往需要先在更清晰、更安全的材料中反复练习的现实需要。它不是把训练做成单线测试，而是让不同能力维度在一个相对连贯的节奏里逐步展开。

\PlaceholderFigure
  {活动路径图占位}
  {建议替换为五类训练活动从“识别”到“表达”的递进关系图。}

\section{从一次训练到一次回看}
系统当前已经具备完整的训练闭环。孩子可以先创建个人档案并选择偏好的机器人伙伴，再进入活动首页开始训练；系统按照当前掌握状态和活动类型发放题目，记录作答结果、反应时间与提示情况；在 Web 场景下，摄像头采集到的图像帧会进入面部情绪识别链路，用于辅助判断儿童是否处于困惑、焦虑或稳定状态；当训练结束后，系统会生成总结信息，并在治疗师看板中呈现情绪正确率、会话记录和掌握度变化。

这一闭环的价值在于，它把原本容易散失的训练细节变成了可被回看的过程信息。对于家长，系统能帮助他们更清楚地理解孩子在一次训练中经历了什么；对于治疗师和机构，系统能帮助他们更高效地复盘某一阶段的训练重点，并逐步积累更稳定的观察基础。

\section{已经做到的，与正在路上的}
当前原型已经实现的核心能力包括：儿童档案创建与管理、五类训练活动入口与主要作答流程、题库驱动的发题机制、会话记录、训练总结、治疗师看板，以及基于摄像头的实时面部情绪识别与规则驱动自适应。这些能力已经足以支撑项目进行演示、试点沟通和后续迭代。

与此同时，后续仍可继续建设的方向也需要清晰区分。语音情绪识别尚未形成完整训练闭环，生理信号接入仍处于预留与规划阶段，更复杂的长期个性化能力也有待进一步验证。因此，本项目现阶段更适合被定义为“具备完整原型闭环的训练支持系统”，而不是已经完成全面多模态落地的成熟产品。

\PlaceholderFigure
  {产品闭环图占位}
  {建议替换为“儿童训练端、状态识别、总结看板”三部分串联示意图。}

\chapter{商业路径与推广思路}

\section{先为谁创造价值}
从价值对象看，星语灵境首先服务的是康复机构、特教支持场景和治疗师团队。对这类用户而言，产品的价值在于帮助其组织标准化训练流程、沉淀过程记录、提高复盘效率，并为家庭沟通提供更直观的依据。对家庭而言，产品的价值则更多体现为降低居家练习的进入门槛，让家长在非专业场景中也能获得相对清晰的训练陪伴工具。

因此，本项目更适合采用“机构优先、家庭延展”的商业路径。先在专业场景中验证可用性与适配度，建立可信案例和合作关系，再逐步把部分能力向家庭场景延伸，会比一开始直接面向大规模消费市场更稳妥。

\section{合作如何发生}
结合当前成熟度，项目早期更适合围绕三类方式展开合作。第一类是机构试点合作，包括试用、共建训练内容、试点部署和教师培训支持；第二类是面向机构的产品服务合作，在原型成熟后逐步过渡到订阅式服务、功能授权或定制化部署；第三类是面向公益项目、课程支持和区域试点的项目制合作，用于扩大真实使用场景并沉淀案例。

在现阶段，是否立刻形成完整收费模型并不是最关键的问题。更关键的是先回答两个问题：产品是否真的能进入训练流程，合作方是否愿意持续使用。只有这两个问题被回答清楚，后续的订阅、授权、部署和家庭延展才有现实基础。

\section{怎样被更多人看见与信任}
推广策略应沿着真实使用关系逐步展开。第一步，是依托合作机构和真实试点场景，形成第一批可展示的使用案例；第二步，是围绕康复机构、特教支持场景和公益合作，建立更稳定的口碑和转介绍机制；第三步，才是在更成熟阶段，逐步向家庭陪伴和更广泛的教育支持场景延伸。

在传播语气上，本项目也需要保持克制。对外表达应避免“治愈”“矫正”“替代专业判断”等夸张说法，而应坚持“理解”“陪伴”“练习”“支持”和“渐进成长”这类更符合项目本质的语言。对 ASD 儿童而言，语言风格本身就是产品态度的一部分。

\chapter{实施计划与试点安排}

\section{先从哪些场景开始}
当前阶段，项目最适合进入的试点对象包括康复机构、特教支持场景以及与 ASD 儿童服务相关的公益项目。与这些场景合作，一方面能够帮助团队在较真实的训练节奏中观察产品是否可用，另一方面也更容易收集来自治疗师、家长和机构管理者的多角色反馈，为后续迭代提供依据。

相比大范围推广，小规模、高质量的试点更符合当前产品成熟度。对于这样一个面向特殊儿童的项目而言，首批合作场景的反馈质量，往往比合作数量更重要。

\section{试点将如何推进}
在实施上，本项目建议按照“场景接入、使用引导、训练开展、结果回看、反馈迭代”五个环节推进。接入阶段重点完成产品部署、角色说明和基本操作培训；使用引导阶段重点帮助儿童熟悉界面、伙伴形象和活动节奏；训练开展阶段观察儿童的参与情况、提示频次和活动完成度；结果回看阶段由家长和治疗师共同阅读训练总结与看板信息；反馈迭代阶段则把实际问题重新带回产品与内容优化中。

这一流程的目标不是尽快把功能全部铺开，而是让每一个环节都能被真实理解、真实使用。对于 ASD 场景而言，能够稳定进入一次训练，往往比快速增加功能更有意义。

\PlaceholderFigure
  {试点流程图占位}
  {建议替换为“接入、训练、回看、反馈”四阶段试点流程图。}

\section{我们重点看什么}
试点验证应重点观察三类指标。第一类是可用性指标，例如儿童是否愿意进入训练、活动完成率如何、支持者是否能顺利理解界面和结果。第二类是适配性指标，例如五类活动是否贴近训练场景、提示方式是否合适、不同活动对不同儿童的接受度是否存在明显差异。第三类是过程性指标，例如情绪正确率变化、提示触发频次、反应时间变化以及支持者对总结信息的实际使用情况。

现阶段更适合承诺的是“具备进入真实试点并积累验证数据的条件”，而不是直接给出疗效性结论。对于本项目而言，先把可用性和适配度做扎实，才有资格进一步讨论更长期的效果评价。

\section{数据与边界的分寸}
由于项目涉及儿童图像、训练记录和情绪信息，数据与伦理要求必须作为实施方案的一部分被明确呈现。进入试点前，需要建立清晰的知情说明与授权流程，明确采集范围、使用目的和保存方式；进入试点后，应坚持最小化采集原则，把状态感知能力定位为训练辅助，而不是替代专业判断的“读心工具”。

这一边界不仅关系到合规，也关系到项目能否被长期信任。面对 ASD 儿童，技术的分寸感本身就是专业性的一部分。

\chapter{资源投入与财务安排}

\section{现阶段最值得投入的地方}
在当前阶段，财务规划的重点不应放在夸大的收入预测上，而应放在资源优先级上。对星语灵境而言，最需要投入的并不是营销噱头，而是产品稳定性、训练内容质量、试点支持能力和合规准备。只有这些基础被做实，项目后续的商业化节奏才有可能真正成立。

因此，资金的主要去向应围绕四件事展开：继续打磨儿童端体验与支持者端反馈能力，完善训练内容与题库组织，完成试点所需的部署与沟通材料，以及补齐隐私说明、授权文本和使用边界等基础工作。

\section{预算如何分配}
考虑到当前项目阶段，预算更适合按结构管理，而非按过于精确的远期报表管理。建议的资源配置重点如下：

\begin{table}[H]
\centering
\small
\begin{tabularx}{0.94\textwidth}{>{\raggedright\arraybackslash}p{0.24\textwidth} >{\centering\arraybackslash}p{0.18\textwidth} X}
\toprule
投入方向 & 建议占比 & 主要用途 \\
\midrule
产品与研发优化 & 40\%--45\% & 儿童端训练流程、支持者看板、状态识别链路与整体稳定性迭代 \\
训练内容与视觉素材 & 15\%--20\% & 题库补充、活动内容完善、插图与界面素材统一、表达风格打磨 \\
试点与用户研究 & 20\%--25\% & 场景部署、使用培训、反馈收集、试点沟通与验证工作 \\
运营与基础部署 & 10\%--15\% & 设备、环境、日常维护与必要的技术支持 \\
合规与文档准备 & 5\%--10\% & 授权说明、隐私材料、展示材料与项目申报文档整理 \\
\bottomrule
\end{tabularx}
\end{table}

\section{资源使用的基本原则}
资源使用应坚持两个原则。第一，优先保障真实使用场景中的关键问题，而不是优先追求表面上的功能丰富度；第二，优先建设可复用的内容和流程，而不是一次性展示型投入。对当前阶段的项目而言，最有价值的投入，是那些能直接提升试点质量、产品稳定性和合作方理解效率的投入。

\chapter{风险、边界与长期信任}

\section{产品是否真正适合孩子}
面向 ASD 儿童的产品，最大的风险并不一定来自“做不出来”，而可能来自“做出来却不合适”。如果互动节奏过快、反馈过强、提示方式过于压迫，孩子即便完成了训练，也不代表训练体验是友好的。因此，产品适配风险始终需要放在首位。

对应的应对策略，是持续围绕低刺激、渐进式、可暂停和可调整的原则打磨产品细节，并在试点中充分收集不同儿童对活动、语言、界面和提示方式的实际反应。

\section{从原型走向使用的挑战}
当前原型已经打通基础闭环，但技术与运营风险仍然存在。例如，面部情绪识别在不同设备、角度和光照条件下的稳定性仍需继续观察；规则驱动自适应虽然可解释，但仍需在更多真实场景中验证其适配度；训练内容和试点支持流程也需要持续优化，避免产品能演示却难持续使用。

这类风险的应对重点，不在于过早承诺更复杂的能力，而在于把当前已经跑通的闭环做稳，把真实场景中的问题一项项消化掉。对于当前阶段而言，稳定比炫目更重要。

\section{隐私与合规底线}
涉及儿童图像、情绪信息和训练记录时，隐私与合规始终是底线要求。项目需要坚持最小化采集、明确授权、有限解释和分角色查看的原则。对于合作机构和家庭，也应明确说明系统的用途边界，即它是训练支持工具，而不是医疗诊断工具，更不是替代专业判断的万能系统。

只有在边界清晰的前提下，项目才有可能获得长期信任。对 ASD 儿童而言，被尊重的边界感，本身就是支持的一部分。

\chapter{社会价值与更长远的意义}

\section{对孩子意味着什么}
对儿童而言，星语灵境的意义不只是多了一套数字化活动，而是多了一个更容易进入训练、也更容易被温和接住的练习环境。孩子可以在更明确、更稳定的节奏中反复练习情绪识别、情境理解和自我表达，这种“允许慢慢来”的体验，本身就是重要价值。

\section{对家庭与机构意味着什么}
对家庭而言，项目能够降低居家练习的启动难度，也能帮助家长更具体地理解孩子在训练中的表现。对机构而言，项目能够把部分标准化训练流程、会话记录和总结反馈组织起来，减少重复记录带来的压力，并为后续沟通提供更清晰的依据。

\section{对行业与公益合作的意义}
在更大的层面上，星语灵境提供了一种“科技向善如何贴近特殊儿童支持场景”的实践路径。它既可以为公益合作、课程支持和康复场景提供更具可操作性的数字工具，也能够为后续相关项目积累关于温和交互、训练记录和多角色协同的经验。

\chapter{团队力量与未来展望}

\section{团队如何各自承担责任}
围绕当前系统建设与后续试点需求，团队需要形成清晰、稳定的职责分工。更合适的方式，不是简单罗列头衔，而是围绕产品落地链条进行分工组织：

\begin{table}[H]
\centering
\small
\begin{tabularx}{0.96\textwidth}{>{\raggedright\arraybackslash}p{0.22\textwidth} X X}
\toprule
职责方向 & 主要任务 & 与当前系统的对应关系 \\
\midrule
项目统筹与产品策划 & 负责整体方案推进、需求梳理、章节叙事、试点沟通与阶段目标管理 & 对应项目定位、产品结构整理、试点推进与对外合作表达 \\
前端与交互实现 & 负责儿童端训练界面、伙伴形象、活动交互、训练总结页与使用体验优化 & 对应当前训练页面、Tino/Tina 伙伴、总结页与主页交互实现 \\
后端与数据能力 & 负责档案管理、会话记录、题库调度、训练报告、看板能力与自适应逻辑维护 & 对应当前档案、答题记录、BKT 状态、训练总结与治疗师看板 \\
内容设计与场景测试 & 负责训练内容组织、活动文案、机构反馈收集、试点执行与迭代建议整理 & 对应五类活动内容打磨、题库完善、机构协同与试点反馈闭环 \\
\bottomrule
\end{tabularx}
\end{table}

这样的分工方式更贴近项目当前阶段。一方面，它能够保证系统继续完善时有人对关键链路负责；另一方面，也能让团队在面对合作机构、家长和评审时，形成清晰一致的表达与执行节奏。

\section{怎样把不同专业放在一起}
除了团队内部职责分工之外，本项目还需要建立外部协作机制，尤其是与特教、康复、心理支持等专业场景的持续连接。对于面向 ASD 儿童的产品来说，真正重要的不是“技术团队单独完成了什么”，而是技术、内容和场景理解能否稳定地合在一起。因此，后续与治疗师、机构顾问和试点合作方的协同，将是产品质量持续提升的重要基础。

\PlaceholderFigure
  {团队协作图占位}
  {建议替换为“团队职责线 + 外部专业协作方”关系示意图。}

\section{未来展望}
未来一段时间内，项目的重点仍然是把现有原型打磨得更稳、更适合进入真实场景。短期内，应继续完善题库内容、提示机制、总结反馈和试点流程，使系统在机构合作中更容易被理解和使用；中期内，可逐步完善角色权限、内容管理与家庭延展能力，让产品从演示原型走向更稳定的服务工具；更长远地看，语音情绪识别、生理信号接入和更细致的个性化支持能力，也可以在真实条件成熟后继续纳入建设路线。

无论产品未来走到哪一步，项目都应保留一个最初的判断：技术的意义，不是替代人与人之间的支持，而是让这种支持变得更细致、更持续，也更容易发生。对于 ASD 儿童而言，真正重要的从来不是系统有多炫目，而是当他们慢慢靠近世界时，是否有一些足够耐心、足够温和的工具与人，愿意一直陪在身边。


\end{document}
